\section{Model and problem statement}\label{sec:model}

\paragraph{Time, capacity and population.}
Time is slotted, $t=1,2,\dots$; we think of a slot as one second and measure rates in Mbit/s.
An access point (a home gateway or a public hotspot) has capacity $C_t>0$ in slot $t$.
The set of active devices changes from slot to slot as devices arrive and leave.
At the start of slot $t$ the operator observes $C_t$ and the data of the currently active devices; it has no information, not even distributional, about later slots.

\paragraph{Tiers.}
Every device belongs to a commercial tier $g\in\{F,P,Y\}$: free, premium, or pay-as-you-go (PAYG).
The tier is the billing identity of the subscriber. It is verified by the operator and cannot be misreported.
Tier $g$ carries a weight $w_g>0$ that expresses how much the operator values a unit of utility delivered to that tier.

\paragraph{Service classes.}
A device declares a service class, which the operator cannot verify.
Class $E$ (elastic) stands for delay-tolerant traffic such as e-mail, backups and downloads.
The real-time (RT) classes $k\in\cK=\{\mathrm{v},\mathrm{g}\}$ are interactive video ($\mathrm{v}$) and online gaming ($\mathrm{g}$).

An RT flow of class $k$ is useful only if it receives a rate in $[\rmin_k,\rmax_k]$ in the current slot, where $0<\rmin_k\le\rmax_k$.
A rate below $\rmin_k$ makes the slot worthless, and rate above $\rmax_k$ brings nothing.
Its per-slot utility is $u_k(z)$, bounded, nondecreasing and concave on $[\rmin_k,\rmax_k]$ (strictly concave if $\rmin_k<\rmax_k$) with $u_k(\rmin_k)>0$.
We use $u_{\mathrm{v}}(z)=1+\log(z/\rmin_{\mathrm{v}})$, which models a call whose quality degrades gracefully down to $\rmin_{\mathrm{v}}$, and a constant $u_{\mathrm{g}}\equiv 1$ at a fixed rate $\rmin_{\mathrm{g}}=\rmax_{\mathrm{g}}=r_{\mathrm{g}}$.
Let
\begin{equation}\label{eq:eta}
\eta_k=\max_{z\in[\rmin_k,\rmax_k]}\frac{u_k(z)}{z},
\end{equation}
the largest utility per unit of rate that a class-$k$ flow can generate; for both utilities above $\eta_k=u_k(\rmin_k)/\rmin_k=1/\rmin_k$.

An elastic device $i$ has a demand cap $d_{i,t}\ge 0$ in slot $t$: the smaller of its link rate and the traffic queued for it at the access point.
For downlink traffic both quantities are measured by the access point itself, so $d_{i,t}$ is \emph{verified}, not declared.
The deficit bound of Section~\ref{sec:envy} holds for arbitrary demand caps; the incentive and optimality results of Sections~\ref{sec:ic} and~\ref{sec:opt} additionally treat the caps as exogenous, as for a device with a standing backlog.
The device receives a rate $x_{i,t}\in[0,d_{i,t}]$.
What an elastic user cares about is how much data it gets, not when: its true preference is an arbitrary nondecreasing function of the cumulative volume $\sum_t x_{i,t}$ received during its sojourn.
The operator scores elastic service with the concave surrogate $w_g\,u(x)$, $u(x)=\log(1+x)$, which yields weighted proportional fairness among elastic devices; the analysis only uses that $u$ is strictly concave, differentiable and strictly increasing, so in principle any $\alpha$-fair utility $((1+x)^{1-\alpha}-1)/(1-\alpha)$ can be substituted, with the corresponding change in the closed form \eqref{eq:xresp} (our implementation and experiments use $\alpha=1$, and we have not checked the hypotheses of Proposition~\ref{prop:reverse} for other values).

\paragraph{Decisions.}
Let $n_{g,k,t}$ be the number of active flows of tier $g$ that declare class $k$.
For every nonempty group $(g,k)$ the operator chooses an admission fraction $a_{g,k,t}\in[0,1]$ and a common rate $z_{g,k,t}\in[\rmin_k,\rmax_k]$; a fraction $a$ of the flows is served at rate $z$ and the others are blocked in this slot.
We write $y_{g,k,t}=a_{g,k,t}z_{g,k,t}$ for the mean rate per flow.
The capacity constraint is
\begin{equation}\label{eq:cap}
\sum_{g,k} n_{g,k,t}\,y_{g,k,t}+\sum_{i\in\cE_t}x_{i,t}\le C_t ,
\end{equation}
where $\cE_t$ is the set of active devices that declare $E$.
The total rate of tier $g$ is $\rho_{g,t}=\sum_k n_{g,k,t}y_{g,k,t}+\sum_{i\in\cE_t:\,g_i=g}x_{i,t}$ and its demand is $\dem_{g,t}=\sum_k n_{g,k,t}\rmax_k+\sum_{i\in\cE_t:\,g_i=g}d_{i,t}$.
The paid tiers come with a \emph{committed rate} $c_g$ per device, the rate the product promises; the committed demand of tier $g$ is $\mathrm{com}_{g,t}=\sum_k n_{g,k,t}\min(\rmax_k,c_g)+\sum_{i\in\cE_t:\,g_i=g}\min(d_{i,t},c_g)$.
The (expected) utility collected in slot $t$ is
\begin{equation}\label{eq:util}
U_t=\sum_{g,k}w_g\,n_{g,k,t}\,a_{g,k,t}\,u_k(z_{g,k,t})+\sum_{i\in\cE_t}w_{g_i}u(x_{i,t}).
\end{equation}

\paragraph{Business constraints.}
The operator's commercial commitments are long-run constraints on tier rates:
\begin{align}
\text{premium service level:}\quad &\limsup_{T\to\infty}\frac1T\sum_{t\le T}\big(\theta_P\,\mathrm{com}_{P,t}-\rho_{P,t}\big)\le 0,\label{eq:ps}\\
\text{PAYG service level:}\quad &\limsup_{T\to\infty}\frac1T\sum_{t\le T}\big(\theta_Y\,\mathrm{com}_{Y,t}-\rho_{Y,t}\big)\le 0,\label{eq:ys}\\
\text{free-tier floor:}\quad &\limsup_{T\to\infty}\frac1T\sum_{t\le T}\big(\min(\phi\,C_t,\dem_{F,t})-\rho_{F,t}\big)\le 0.\label{eq:ff}
\end{align}
Constraint~\eqref{eq:ps} says that the premium tier receives on average at least a fraction $\theta_P$ of its committed demand.
PAYG subscribers pay a price $p$ per unit of volume, so the revenue rate is $p\,\rho_{Y,t}$ and \eqref{eq:ys} guarantees that the operator captures at least a fraction $\theta_Y$ of the PAYG revenue that the committed rates put on offer.
Constraint~\eqref{eq:ff} reserves a share $\phi$ of the capacity for the free tier whenever free users want it, as venues that advertise free access typically must.
We write $e_{j,t}$, $j\in\{PS,YS,FF\}$, for the three summands, so that each constraint reads ``the time average of $e_j$ is nonpositive''; we take $\theta_P,\theta_Y,\phi\in[0,1]$.
All three are linear in the tier rates. This matters for the analysis: the price that such a constraint puts on a unit of rate is the same for every class inside a tier.

\paragraph{Strategic behaviour and class envy.}
Protection of RT flows is valuable, so an elastic device may be tempted to declare an RT class.
If it does, it receives the RT offer of its tier; in a slot in which the group $(g,k)$ is offered $(a,z)$, the expected volume it can actually use is
\begin{equation}\label{eq:ref}
r_{i,k,t}=a_{g_i,k,t}\,\min\big(z_{g_i,k,t},\,d_{i,t}\big),
\end{equation}
with $r_{i,k,t}=0$ when the group is empty.
We call $r_{i,k,t}$ the \emph{reference} of device $i$ for class $k$.

\begin{definition}[Class envy]\label{def:envy}
The class envy of elastic device $i$ towards class $k$ over a window of slots $[s,t]$ inside its sojourn is $\sum_{q=s}^{t}\big(r_{i,k,q}-x_{i,q}\big)$.
An allocation rule is \emph{$B$-class-envy-free} if this quantity is at most $B$ for every device, class and window, on every sample path.
\end{definition}

Class envy is measured in volume because volume is what an elastic user values.
A rule with bounded class envy may still serve RT flows first when capacity is short; it only has to return the difference later.
A rule with $B=0$ can never favour an RT flow over an elastic tier-mate in any slot, so it offers no protection at all.

\paragraph{Problem.}
Design an online rule that (i)~maximises the long-run average of $U_t$, (ii)~satisfies \eqref{eq:ps}--\eqref{eq:ff}, and (iii)~is $B$-class-envy-free for a small $B$, without knowledge of future capacities, arrivals, departures or demands.
Section~\ref{sec:analysis} makes precise the benchmark against which (i) is measured and the sense in which (iii) limits the gain from misdeclaring the class.
